\Needspace{12\baselineskip}
\originalchapter{水溶液中的静电作用}
\originalsection{Bjerrum长度和屏蔽}
生物分子常带电，但在水中的电荷作用不能按真空库仑力估计。采用SI单位，两个单位电荷的能量为 $e^2/(4\pi\epsilon r)$；定义 Bjerrum 长度 $\ell_B=e^2/(4\pi\epsilon\kb T)$，使这一能量等于 $\kb T$。讲义源文有时把 $4\pi$ 吸收到介电常数中，本册用 $\ell_B$ 表示结果以免混淆。水中室温典型值约 $0.7\,\mathrm{nm}$，而它不是所有溶剂通用的常数。

对价态 $z_\alpha$、体相密度 $n_\alpha^0$ 的可移动离子，平均场局部平衡给出 $n_\alpha(\mathbf r)=n_\alpha^0\e^{-\beta z_\alpha e\phi}$。与 Poisson 方程结合便得到 Poisson--Boltzmann 方程。若 $|\beta z_\alpha e\phi|\ll1$，在电中性的体相附近线性化有
\[
(\nabla^2-\kappa_D^2)\phi=-\frac{\rho_{\rm fixed}}\epsilon,
\qquad \kappa_D^2=\frac{e^2}{\epsilon\kb T}\sum_\alpha z_\alpha^2n_\alpha^0.
\]
点电荷的球对称解为 $\phi=q\e^{-\kappa_D r}/(4\pi\epsilon r)$。盐越多，屏蔽长度 $\lambda_D=\kappa_D^{-1}$ 越短。接近高电荷表面、强关联多价离子或分子尺度时，线性化和连续平均场都可能失效。

\originalsection{无盐平面双电层}
取负表面电荷密度 $-\sigma$（$\sigma>0$），溶液在 $y>0$，只有价态 $z>0$ 的反离子。定义 $u=\beta ze[\phi(y)-\phi(0)]$。Boltzmann分布 $n=n_0\e^{-u}$ 和 Poisson 方程给 $u''=-\beta z^2e^2 n_0\e^{-u}/\epsilon$。满足远处电场趋零的解是
\[
u(y)=2\ln(1+y/\mu),\qquad
n(y)=\frac1{2\pi z^2\ell_B(y+\mu)^2},\qquad
\mu=\frac{e}{2\pi z\ell_B\sigma}.
\]
边界 $E(0)=-\sigma/\epsilon$ 确定 $\mu$，其中 $E=-\phi'$ 指向负膜。积分 $ze\int_0^\infty n\dd y=\sigma$ 核查整体电中性。这种密度按幂律衰减，和有盐的指数屏蔽不同；将它随意用在含体相盐的溶液中是不一致的。

\originalsection{线电荷与反离子凝聚}
离子位于半径 $a<r<R$ 的柱形区域，线电荷吸引产生能量 $2\zeta\kb T\ln(r/a)$。径向面积元 $2\pi r\dd r$ 与 Boltzmann 权重共同给出径向概率 $p_r(r)\propto r^{1-2\zeta}$。因此决定凝聚的是能量与二维径向熵的竞争，而非仅仅“势能在无穷远增加”。归一化积分在 $R\to\infty$ 时当且仅当 $\zeta>1$ 有限。

写线电荷数密度为 $n$，则 $\zeta=zn\ell_B$。Manning图像假定部分反离子凝聚，使剩余有效线电荷达到 $zn_{\rm eff}\ell_B=1$，故凝聚电荷比例为 $1-1/\zeta$。它是长、均匀、细线电荷的近似；有限长度、离子关联和有限盐浓度需要更细模型。其有限容器推导和径向矩的不同阈值在所选原题解答中给出。

\originalsection{电离平衡与介电环境}
带电基团是否电离由化学自由能和溶剂决定。最简反应 $AH\rightleftarrows A^-+H^+$ 的平衡常数 $K_a=a_{A^-}a_{H^+}/a_{AH}$ 使用无量纲活度。稀溶液近似活度等于浓度与标准浓度之比，可得 $[A^-]/[AH]=10^{\mathrm{pH}-\mathrm{p}K_a}$。表面有电势时，质子局部浓度还带Boltzmann因子；因而把真空中离开一对电荷的库仑能当作水中酸度并不成立。

水的高介电常数降低电荷能，但连续介质在离子半径内不能外推到 $r=0$。一个半径 $a$ 的离子放入介电介质，充电模型给自能尺度 $q^2/(8\pi\epsilon a)$。更换溶剂的差值能说明离子进入水相的偏好，但半径、局部溶剂结构和化学键能都影响精确值。本节重建源讲义电离讨论的物理链，不使用来源中的第三方分子图片。

\originalsection{Poisson--Boltzmann的自由能推导}
把固定电荷密度 $\rho_f$ 与移动离子密度 $n_\alpha$ 合并。电场能为 $\tfrac12\int\rho\phi\dd^3r$，理想混合熵项为 $\kb T\sum_\alpha\int n_\alpha[\ln(n_\alpha/n_\alpha^0)-1]\dd^3r$。对 $n_\alpha$ 变分时电场能的变化为 $z_\alpha e\phi\,\delta n_\alpha$，所以驻点条件给 $n_\alpha=n_\alpha^0\e^{-\beta z_\alpha e\phi}$。再代入 $-\epsilon\nabla^2\phi=\rho_f+\sum z_\alpha en_\alpha$ 才得到PB方程。离子相关在此平均场自由能中没有被计算。

对称1:1盐密度各为 $n_s$，令 $\psi=\beta e\phi$，平面区方程为 $\psi''=\kappa_D^2\sinh\psi$。乘 $\psi'$ 并利用无穷远 $\psi=\psi'=0$，得 $(\psi')^2=2\kappa_D^2(\cosh\psi-1)$。一阶积分的解可写成
\[
\tanh\frac{\psi(y)}4=\tanh\frac{\psi(0)}4\,\e^{-\kappa_Dy}.
\]
远场总恢复指数衰减，但近表面可能不线性。带符号表面电荷 $\sigma_s$ 与边界 $-\epsilon\phi'(0)=\sigma_s$ 给Grahame关系
\[
\sigma_s^2=8\epsilon\kb T n_s\sinh^2\frac{\psi(0)}2.
\]
它说明高表面电势并不与表面电荷保持简单线性关系。

\originalsection{无盐平面解的积分与压力核对}
无盐方程中 $u''=-C\e^{-u}$。乘 $u'$ 后积分，远处电场为零使 $(u')^2=2C\e^{-u}$。取 $u(0)=0$、向溶液方向 $u'>0$，积分为 $\e^{u/2}=1+\sqrt{C/2}\,y$，故 $\mu=\sqrt{2/C}$。这样既得到正文密度，也看到 $n_0$ 不是可任意指定的体相密度，它由表面电荷和整体中性共同确定。

平面体系法向总压力 $P=\kb T\sum n_\alpha-\epsilon E^2/2$ 为常数。单面无盐远处 $P=0$，表面接触密度满足 $\kb T n(0)=\sigma^2/(2\epsilon)$。两相同带电膜之间中点 $E=0$，所以排斥压力等于中点渗透压。这是检查原题两板解的独立方法，且不需要把两个孤立单板解直接叠加；PB是非线性方程，叠加一般不成立。

\originalsection{有限容器与凝聚极限的次序}
柱形离子的径向归一化积分 $I=(R^{2-2\zeta}-a^{2-2\zeta})/(2-2\zeta)$，在 $\zeta=1$ 应取极限 $I=\ln(R/a)$。所谓凝聚过渡是先讨论 $R/a\to\infty$ 的分布：$\zeta>1$ 时有限半径内留住非零概率；$\zeta\le1$ 时概率向容器尺度扩散。但平均半径的分子是 $\int r^{2-2\zeta}\dd r$，要到 $\zeta>3/2$ 才收敛。因此“已经凝聚”与“平均半径有限”是不同阈值。将每个矩的阈值分别计算，比只凭一条曲线认定相态可靠。

裸DNA电荷间距 $b$ 给 $\zeta=z\ell_B/b$。在Manning有效电荷近似下，剩余电荷比例 $1/\zeta$；对于多价离子这只是长线平均场的估计。多价离子的局部桥接、强关联甚至电荷反转都超出本文模型，不能用简单凝聚比例代替这些机制。
